\section{Тема, мета та вихідні дані}

\subsection*{Тема КР}

Рій безпілотних літальних апаратів (БПЛА) для автономної інспекції складського приміщення з потоковою IoT-телеметрією у хмарний стек та алгоритмом естафетної передачі місії при досягненні критичного рівня заряду акумулятора.

\subsection*{Мета КР}

Спроєктувати, реалізувати та продемонструвати в симуляційному середовищі систему \projname{}, яка забезпечує:

\begin{enumerate}[itemsep=2pt]
  \item автономне обстеження складського приміщення роєм 3 БПЛА з призначенням кожному дрону окремого коридора між стелажами;
  \item точне слідування квінтичними поліноміальними траєкторіями \cite{Lynch2017} та геометричний контроль орієнтації на SE(3) \cite{Lee2010,Mellinger2011};
  \item планування маршруту покриття площі коридора методом \emph{Boustrophedon Cellular Decomposition} (BCD) \cite{Choset2000};
  \item потокову передачу телеметрії (позиція, швидкість, заряд батареї, прогрес покриття) за протоколом MQTT v5 \cite{MQTT5} з подальшою візуалізацією у дашборді Grafana;
  \item автоматичну естафетну передачу місії між двома дронами: при падінні заряду батареї нижче порогового значення другий дрон перехоплює \emph{поточні світові координати} першого і продовжує сканування.
\end{enumerate}

\subsection*{Вихідні дані}

\begin{itemize}[itemsep=2pt]
  \item Симуляційне середовище: NVIDIA Isaac Sim 4.0+ з надбудовою Pegasus Simulator \cite{Pegasus2024}, сцена «Full Warehouse» з NVIDIA Omniverse Asset Store.
  \item Апаратна платформа дрона (в симуляції): квадрокоптер Iris з готовою фізичною моделлю Pegasus, маса $m = 1{,}5$~кг.
  \item IoT-стек: Eclipse Mosquitto~2.x, Telegraf~1.30, InfluxDB~2.7, Grafana~10.4; оркестрація через Docker Compose.
  \item Мова реалізації: Python 3.11; бібліотеки \texttt{paho-mqtt}, \texttt{pydantic}, \texttt{numpy}.
  \item Топологія об'єкту: три паралельних коридори між стелажами, вісь Y, довжина 11~м (від $y=-5{,}5$ до $y=+5{,}5$~м), висота сканування $z=3{,}0$~м.
  \item Цільова швидкість крейсерського польоту: $v_{\text{cruise}} = 0{,}5$~м/с.
\end{itemize}
